\documentclass[11pt]{amsart}
\usepackage[margin=1.05in]{geometry}
\usepackage{amsmath,amssymb}
\usepackage{graphicx}
\usepackage[colorlinks=true,linkcolor=blue,citecolor=blue,urlcolor=blue]{hyperref}

\title{Dungey's Averaged \(A_\alpha\) Ritt-Probability Conjecture}
\author{fa\_banach\_001 literature-status packet}
\date{2026-07-03}

\begin{document}
\maketitle

\section*{Status}

This is a literature-already-answered packet.  The source paper is Nick
Dungey's \emph{Subordinated discrete semigroups of operators},
arXiv:0801.4557.  Dungey defines, for \(\varepsilon\in(0,1]\),
\[
  B=\varepsilon^{-1}\int_0^\varepsilon A_\alpha\,d\alpha
\]
and conjectures on PDF page 19 that \(B\) belongs to the class \(\mathcal A\)
of Ritt probabilities.

\begin{center}
\includegraphics[width=0.9\linewidth]{figures/open_problem_crop.png}
\end{center}

\section*{Answer Source}

Alexander Gomilko and Yuri Tomilov, \emph{On discrete subordination of power
bounded and Ritt operators}, arXiv:1504.00563, explicitly identify this as a
question left open by Dungey.  In their notation,
\[
  h_\varepsilon(\lambda)
  =1-\frac1\varepsilon\int_0^\varepsilon (1-\lambda)^\alpha\,d\alpha
  =1-\frac{(1-\lambda)^\varepsilon-1}{\varepsilon\log(1-\lambda)} .
\]
Theorem 7.1 proves a geometric criterion for regular Hausdorff functions:
if \(1-h(\mathbb D)\) lies in a sector \(\overline{\Sigma}_\gamma\), then
\(h(T)\) is Ritt for every power-bounded operator \(T\).  Example 7.4 verifies
this criterion for \(h_\varepsilon\), and concludes that \(h_\varepsilon(T)\)
is Ritt of angle \(\varepsilon\pi/2\) for \(0<\varepsilon<1\).  The paragraph
immediately following also covers the endpoint \(h_1\).

\begin{center}
\includegraphics[width=0.9\linewidth]{figures/supporting_answer_crop.png}
\end{center}

\section*{Identification}

Dungey's probability \(B\) has generating function
\[
  \phi_B(\lambda)=h_\varepsilon(\lambda).
\]
Dungey's Theorem 3.1 characterizes \(B\in\mathcal A\) by the equivalent
operator statement that \(\Psi(B;T)=h_\varepsilon(T)\) is a Ritt operator for
every power-bounded \(T\).  Gomilko--Tomilov prove exactly this statement and
explicitly state that it settles Dungey's conjecture.

\section*{References}

\begin{thebibliography}{9}
\bibitem{Dungey}
N. Dungey, \emph{Subordinated discrete semigroups of operators},
arXiv:0801.4557.

\bibitem{GT}
A. Gomilko and Y. Tomilov, \emph{On discrete subordination of power bounded
and Ritt operators}, arXiv:1504.00563.
\end{thebibliography}

\end{document}
